\section{AHSME 1960}
        \begin{enumerate}
        	\item (\textit{AHSME 1960}) If $2$ is a solution (root) of $x^3+hx+10=0$, then $h$ equals:


 $\textbf{(A)}10\qquad \textbf{(B)}9 \qquad \textbf{(C)}2\qquad \textbf{(D)}-2\qquad \textbf{(E)}-9$


 

	\item (\textit{AHSME 1960}) It takes $5$ seconds for a clock to strike $6$ o'clock beginning at $6{:}00$ o'clock precisely. 
If the strikings are uniformly spaced, how long, in seconds, does it take to strike $12$ o'clock?


 $\textbf{(A)}9\frac{1}{5}\qquad \textbf{(B )}10\qquad \textbf{(C )}11\qquad \textbf{(D )}14\frac{2}{5}\qquad \textbf{(E )}\text{none of these}$


 

	\item (\textit{AHSME 1960}) Applied to a bill for $\textdollar{10,000}$ the difference between a discount of $40$\% and two successive discounts of $36$\% and $4$\%, 
expressed in dollars, is:


 $\textbf{(A) }0\qquad \textbf{(B) }144\qquad \textbf{(C) }256\qquad \textbf{(D) }400\qquad \textbf{(E) }416$


 

	\item (\textit{AHSME 1960}) Each of two angles of a triangle is $60^{\circ}$ and the included side is $4$ inches. The area of the triangle, in square inches, is:


 $\textbf{(A) }8\sqrt{3}\qquad \textbf{(B) }8\qquad \textbf{(C) }4\sqrt{3}\qquad \textbf{(D) }4\qquad \textbf{(E) }2\sqrt{3}$


 

	\item (\textit{AHSME 1960}) The number of distinct points common to the graphs of $x^2+y^2=9$ and $y^2=9$ is:


 $\textbf{(A) }\text{infinitely many}\qquad \textbf{(B) }\text{four}\qquad \textbf{(C) }\text{two}\qquad \textbf{(D) }\text{one}\qquad \textbf{(E) }\text{none}$


 

	\item (\textit{AHSME 1960}) The circumference of a circle is $100$ inches. The side of a square inscribed in this circle, expressed in inches, is:


 $\textbf{(A) }\frac{25\sqrt{2}}{\pi}\qquad \textbf{(B) }\frac{50\sqrt{2}}{\pi}\qquad \textbf{(C) }\frac{100}{\pi}\qquad \textbf{(D) }\frac{100\sqrt{2}}{\pi}\qquad \textbf{(E) }50\sqrt{2}$


 

	\item (\textit{AHSME 1960}) Circle $I$ passes through the center of, and is tangent to, circle $II$. The area of circle $I$ is $4$ square inches. 
Then the area of circle $II$, in square inches, is:


 $\textbf{(A) }8\qquad \textbf{(B) }8\sqrt{2}\qquad \textbf{(C) }8\sqrt{\pi}\qquad \textbf{(D) }16\qquad \textbf{(E) }16\sqrt{2}$


 

	\item (\textit{AHSME 1960}) The number $2.5252525\ldots$ can be written as a fraction. 
When reduced to lowest terms the sum of the numerator and denominator of this fraction is:


 $\textbf{(A) }7\qquad \textbf{(B) } 29\qquad \textbf{(C) }141\qquad \textbf{(D) }349\qquad \textbf{(E) }\text{none of these}$


 

	\item (\textit{AHSME 1960}) The fraction $\frac{a^2+b^2-c^2+2ab}{a^2+c^2-b^2+2ac}$ is (with suitable restrictions of the values of a, b, and c):


 $\text{(A) irreducible}\qquad$


 $\text{(B) reducible to negative 1}\qquad$


 $\text{(C) reducible to a polynomial of three terms}\qquad$


 $\text{(D) reducible to} \frac{a-b+c}{a+b-c}\qquad$


 $\text{(E) reducible to} \frac{a+b-c}{a-b+c}$


 

	\item (\textit{AHSME 1960}) Given the following six statements:

 \[\text{(1) All women are good drivers}\]

 \[\text{(2) Some women are good drivers}\]

 \[\text{(3) No men are good drivers}\]

 \[\text{(4) All men are bad drivers}\]

 \[\text{(5) At least one man is a bad driver}\]

 \[\text{(6) All men are good drivers.}\]


 

The statement that negates statement $(6)$ is:


 

$\textbf{(A) }(1)\qquad \textbf{(B) }(2)\qquad \textbf{(C) }(3)\qquad \textbf{(D) }(4)\qquad \textbf{(E) }(5)$


 

	\item (\textit{AHSME 1960}) For a given value of $k$ the product of the roots of $x^2-3kx+2k^2-1=0$
is $7$. The roots may be characterized as:


 $\textbf{(A) }\text{integral and positive} \qquad\textbf{(B) }\text{integral and negative} \qquad \\ \textbf{(C) }\text{rational, but not integral} \qquad\textbf{(D) }\text{irrational} \qquad\textbf{(E) } \text{imaginary}$


 

	\item (\textit{AHSME 1960}) The locus of the centers of all circles of given radius $a$, in the same plane, passing through a fixed point, is:


 $\textbf{(A) }\text{a point}\qquad \textbf{(B) }\text{ a straight line}\qquad \textbf{(C) }\text{two straight lines}\qquad \textbf{(D) }\text{a circle}\qquad  \textbf{(E) }\text{two circles}$


 

	\item (\textit{AHSME 1960}) The polygon(s) formed by $y=3x+2, y=-3x+2$, and $y=-2$, is (are):


 $\textbf{(A) }\text{an equilateral triangle}\qquad\textbf{(B) }\text{an isosceles triangle} \qquad\textbf{(C) }\text{a right triangle} \qquad \\ \textbf{(D) }\text{a triangle and a trapezoid}\qquad\textbf{(E) }\text{a quadrilateral}$


 

	\item (\textit{AHSME 1960}) If $a$ and $b$ are real numbers, the equation $3x-5+a=bx+1$ has a unique solution $x$ [The symbol $a \neq 0$ means that $a$ is different from zero]:


 $\text{(A) for all a and b} \qquad \text{(B) if a }\neq\text{2b}\qquad \text{(C) if a }\neq 6\qquad \\ \text{(D) if b }\neq 0\qquad \text{(E) if b }\neq 3$


 

	\item (\textit{AHSME 1960}) Triangle $I$ is equilateral with side $A$, perimeter $P$, area $K$, and circumradius $R$ (radius of the circumscribed circle). 
Triangle $II$ is equilateral with side $a$, perimeter $p$, area $k$, and circumradius $r$. If $A$ is different from $a$, then: 


 $\textbf{(A)}\ P:p = R:r \text{ } \text{only sometimes} \qquad \textbf{(B)}\ P:p = R:r \text{ } \text{always}\qquad \\ \textbf{(C)}\ P:p = K:k \text{ } \text{only sometimes} \qquad \textbf{(D)}\ R:r = K:k \text{ } \text{always}\qquad \textbf{(E)}\ R:r = K:k \text{ } \text{only sometimes}$


 

	\item (\textit{AHSME 1960}) In the numeration system with base $5$, counting is as follows: $1, 2, 3, 4, 10, 11, 12, 13, 14, 20,\ldots$.
The number whose description in the decimal system is $69$, when described in the base $5$ system, is a number with:


 $\textbf{(A)}\ \text{two consecutive digits} \qquad\textbf{(B)}\ \text{two non-consecutive digits} \qquad \\ \textbf{(C)}\ \text{three consecutive digits} \qquad\textbf{(D)}\ \text{three non-consecutive digits} \qquad \\ \textbf{(E)}\ \text{four digits}$


 

	\item (\textit{AHSME 1960}) The formula $N=8 \times 10^{8} \times x^{-3/2}$ gives, for a certain group, the number of individuals whose income exceeds $x$ dollars. 
The lowest income, in dollars, of the wealthiest $800$ individuals is at least:


 $\textbf{(A)}\ 10^4\qquad \textbf{(B)}\ 10^6\qquad \textbf{(C)}\ 10^8\qquad \textbf{(D)}\ 10^{12} \qquad \textbf{(E)}\ 10^{16}$


 

	\item (\textit{AHSME 1960}) The pair of equations $3^{x+y}=81$ and $81^{x-y}=3$ has:


 $\textbf{(A)}\ \text{no common solution} \qquad \\ \textbf{(B)}\ \text{the solution} \text{ } x=2, y=2\qquad \\ \textbf{(C)}\ \text{the solution} \text{ } x=2\frac{1}{2}, y=1\frac{1}{2} \qquad \\ \textbf{(D)}\text{ a common solution in positive and negative integers} \qquad \\ \textbf{(E)}\ \text{none of these}$


 

	\item (\textit{AHSME 1960}) Consider equation $I: x+y+z=46$ where $x, y$, and $z$ are positive integers, and equation $II: x+y+z+w=46$, 
where $x, y, z$, and $w$ are positive integers. Then


 $\textbf{(A)}\ \text{I can be solved in consecutive integers} \qquad \\ \textbf{(B)}\ \text{I can be solved in consecutive even integers} \qquad \\ \textbf{(C)}\ \text{II can be solved in consecutive integers} \qquad \\ \textbf{(D)}\ \text{II can be solved in consecutive even integers} \qquad \\ \textbf{(E)}\ \text{II can be solved in consecutive odd integers}$


 

	\item (\textit{AHSME 1960}) The coefficient of $x^7$ in the expansion of $(\frac{x^2}{2}-\frac{2}{x})^8$ is:


 $\textbf{(A)}\ 56\qquad \textbf{(B)}\ -56\qquad \textbf{(C)}\ 14\qquad \textbf{(D)}\ -14\qquad \textbf{(E)}\ 0$


 

	\item (\textit{AHSME 1960}) The diagonal of square $I$ is $a+b$. The perimeter of square $II$ with twice the area of $I$ is:


 $\textbf{(A)}\ (a+b)^2\qquad \textbf{(B)}\ \sqrt{2}(a+b)^2\qquad \textbf{(C)}\ 2(a+b)\qquad \textbf{(D)}\ \sqrt{8}(a+b) \qquad \textbf{(E)}\ 4(a+b)$


 

	\item (\textit{AHSME 1960}) The equality $(x+m)^2-(x+n)^2=(m-n)^2$, where $m$ and $n$ are unequal non-zero constants, is satisfied by $x=am+bn$, where:


 $\textbf{(A)}\ a = 0, b \text{ } \text{has a unique non-zero value}\qquad \\ \textbf{(B)}\ a = 0, b \text{ } \text{has two non-zero values}\qquad \\ \textbf{(C)}\ b = 0, a \text{ } \text{has a unique non-zero value}\qquad \\ \textbf{(D)}\ b = 0, a \text{ } \text{has two non-zero values}\qquad \\ \textbf{(E)}\ a \text{ } \text{and} \text{ } b \text{ } \text{each have a unique non-zero value}$


 

	\item (\textit{AHSME 1960}) The radius $R$ of a cylindrical box is $8$ inches, the height $H$ is $3$ inches. 
The volume $V = \pi R^2H$ is to be increased by the same fixed positive amount when $R$ 
is increased by $x$ inches as when $H$ is increased by $x$ inches. This condition is satisfied by:


 $\textbf{(A)}\ \text{no real value of x} \qquad \\ \textbf{(B)}\ \text{one integral value of x} \qquad \\ \textbf{(C)}\ \text{one rational, but not integral, value of x} \qquad \\ \textbf{(D)}\ \text{one irrational value of x}\qquad \\ \textbf{(E)}\ \text{two real values of x}$


 

	\item (\textit{AHSME 1960}) If $\log_{2x}216 = x$, where $x$ is real, then $x$ is:


 $\textbf{(A)}\ \text{A non-square, non-cube integer}\qquad$
$\textbf{(B)}\ \text{A non-square, non-cube, non-integral rational number}\qquad$
$\textbf{(C)}\ \text{An irrational number}\qquad$
$\textbf{(D)}\ \text{A perfect square}\qquad$
$\textbf{(E)}\ \text{A perfect cube}$


 

	\item (\textit{AHSME 1960}) Let $m$ and $n$ be any two odd numbers, with $n$ less than $m$. 
The largest integer which divides all possible numbers of the form $m^2-n^2$ is:


 $\textbf{(A)}\ 2\qquad \textbf{(B)}\ 4\qquad \textbf{(C)}\ 6\qquad \textbf{(D)}\ 8\qquad \textbf{(E)}\ 16$


 

	\item (\textit{AHSME 1960}) Find the set of $x$-values satisfying the inequality $\left|\frac{5-x}{3}\right| < 2$. [The symbol $|a|$ means $+a$ if $a$ is positive, 
$-a$ if $a$ is negative, $0$ if $a$ is zero. The notation $1 < a < 2$ means that $a$ can have any value between $1$ and $2$, excluding $1$ and $2$.]


 $\textbf{(A)}\ 1 < x < 11\qquad \textbf{(B)}\ -1 < x < 11\qquad \textbf{(C)}\ x < 11\qquad \textbf{(D)}\ x > 11\qquad \textbf{(E)}\ |x| < 6$


 

	\item (\textit{AHSME 1960}) Let $S$ be the sum of the interior angles of a polygon $P$ for which each interior angle is $7\frac{1}{2}$ times the 
exterior angle at the same vertex. Then


 $\textbf{(A)}\ S=2660^{\circ} \text{ } \text{and} \text{ } P \text{ } \text{may be regular}\qquad \\ \textbf{(B)}\ S=2660^{\circ} \text{ } \text{and} \text{ } P \text{ } \text{is not regular}\qquad \\ \textbf{(C)}\ S=2700^{\circ} \text{ } \text{and} \text{ } P \text{ } \text{is regular}\qquad \\ \textbf{(D)}\ S=2700^{\circ} \text{ } \text{and} \text{ } P \text{ } \text{is not regular}\qquad \\ \textbf{(E)}\ S=2700^{\circ} \text{ } \text{and} \text{ } P \text{ } \text{may or may not be regular}$


 

	\item (\textit{AHSME 1960}) The equation $x-\frac{7}{x-3}=3-\frac{7}{x-3}$ has:


 $\textbf{(A)}\ \text{infinitely many integral roots}\qquad\textbf{(B)}\ \text{no root}\qquad\textbf{(C)}\ \text{one integral root}\qquad$
$\textbf{(D)}\ \text{two equal integral roots}\qquad\textbf{(E)}\ \text{two equal non-integral roots}$


 

	\item (\textit{AHSME 1960}) Five times $A$'s money added to $B$'s money is more than $\$51.00$. Three times $A$'s money minus $B$'s money is $\$21.00$.
If $a$ represents $A$'s money in dollars and $b$ represents $B$'s money in dollars, then:


 $\textbf{(A)}\ a>9, b>6 \qquad \textbf{(B)}\ a>9, b<6 \qquad \textbf{(C)}\ a>9, b=6\qquad \textbf{(D)}\ a>9, \text{but we can put no bounds on} \text{ } b\qquad \textbf{(E)}\ 2a=3b$


 

	\item (\textit{AHSME 1960}) Given the line $3x+5y=15$ and a point on this line equidistant from the coordinate axes. Such a point exists in:


 $\textbf{(A)}\ \text{none of the quadrants}\qquad\textbf{(B)}\ \text{quadrant I only}\qquad\textbf{(C)}\ \text{quadrants I, II only}\qquad$
$\textbf{(D)}\ \text{quadrants I, II, III only}\qquad\textbf{(E)}\ \text{each of the quadrants}$


 

	\item (\textit{AHSME 1960}) For $x^2+2x+5$ to be a factor of $x^4+px^2+q$, the values of $p$ and $q$ must be, respectively:


 $\textbf{(A)}\ -2, 5\qquad \textbf{(B)}\ 5, 25\qquad \textbf{(C)}\ 10, 20\qquad \textbf{(D)}\ 6, 25\qquad \textbf{(E)}\ 14, 25$


 

	\item (\textit{AHSME 1960}) In this figure the center of the circle is $O$. $AB \perp BC$, $ADOE$ is a straight line, $AP = AD$, and $AB$ has a length twice the radius. Then:


 
                    \begin{figure}[H]
                        \centering
                        \includegraphics[width=0.5\textwidth]{Problem/AHSME_I_1960_32_Problem_Figure_1.png}
                    \end{figure}
                    


 $\textbf{(A)} AP^2 = PB \times AB\qquad \\ \textbf{(B)}\ AP \times DO = PB \times AD\qquad \\ \textbf{(C)}\ AB^2 = AD \times DE\qquad \\ \textbf{(D)}\ AB \times AD = OB \times AO\qquad \\ \textbf{(E)}\ \text{none of these}$


 

	\item (\textit{AHSME 1960}) You are given a sequence of $58$ terms; each term has the form $P+n$ where $P$ stands for the product $2 \times 3 \times 5 \times\ldots \times 61$
of all prime numbers less than or equal to $61$, and $n$ takes, successively, the values $2, 3, 4,\ldots, 59$. 
Let $N$ be the number of primes appearing in this sequence. Then $N$ is:


 $\textbf{(A)}\ 0\qquad \textbf{(B)}\ 16\qquad \textbf{(C)}\ 17\qquad \textbf{(D)}\ 57\qquad \textbf{(E)}\ 58$


 

	\item (\textit{AHSME 1960}) Two swimmers, at opposite ends of a $90$-foot pool, start to swim the length of the pool, 
one at the rate of $3$ feet per second, the other at $2$ feet per second. 
They swim back and forth for $12$ minutes. Allowing no loss of times at the turns, find the number of times they pass each other. 


 $\textbf{(A)}\ 24\qquad \textbf{(B)}\ 21\qquad \textbf{(C)}\ 20\qquad \textbf{(D)}\ 19\qquad \textbf{(E)}\ 18$


 

	\item (\textit{AHSME 1960}) From point $P$ outside a circle, with a circumference of $10$ units, a tangent is drawn. 
Also from $P$ a secant is drawn dividing the circle into unequal arcs with lengths $m$ and $n$. 
It is found that $t_1$, the length of the tangent, is the mean proportional between $m$ and $n$. 
If $m$ and $t$ are integers, then $t$ may have the following number of values:


 $\textbf{(A)}\ \text{zero}\qquad \textbf{(B)}\ \text{one}\qquad \textbf{(C)}\ \text{two}\qquad \textbf{(D)}\ \text{three}\qquad \textbf{(E)}\ \text{infinitely many}$


 

	\item (\textit{AHSME 1960}) Let $s_1, s_2, s_3$ be the respective sums of $n, 2n, 3n$ terms of the same arithmetic progression with $a$ as the first term and $d$ 
as the common difference. Let $R=s_3-s_2-s_1$. Then $R$ is dependent on:


 $\textbf{(A)}\ a\text{ }\text{and}\text{ }d\qquad \textbf{(B)}\ d\text{ }\text{and}\text{ }n\qquad \textbf{(C)}\ a\text{ }\text{and}\text{ }n\qquad \textbf{(D)}\ a, d,\text{ }\text{and}\text{ }n\qquad  \textbf{(E)}\ \text{neither} \text{ } a \text{ } \text{nor} \text{ } d \text{ } \text{nor} \text{ } n$


 

	\item (\textit{AHSME 1960}) The base of a triangle is of length $b$, and the altitude is of length $h$. 
A rectangle of height $x$ is inscribed in the triangle with the base of the rectangle in the base of the triangle. The area of the rectangle is:


 $\textbf{(A)}\ \frac{bx}{h}(h-x)\qquad \textbf{(B)}\ \frac{hx}{b}(b-x)\qquad \textbf{(C)}\ \frac{bx}{h}(h-2x)\qquad \textbf{(D)}\ x(b-x)\qquad \textbf{(E)}\ x(h-x)$


 

	\item (\textit{AHSME 1960}) In this diagram $AB$ and $AC$ are the equal sides of an isosceles $\triangle ABC$, in which is inscribed equilateral $\triangle DEF$. 
Designate $\angle BFD$ by $a$, $\angle ADE$ by $b$, and $\angle FEC$ by $c$. Then:


 
                    \begin{figure}[H]
                        \centering
                        \includegraphics[width=0.5\textwidth]{Problem/AHSME_I_1960_38_Problem_Figure_1.png}
                    \end{figure}
                    


 $\textbf{(A)}\ b=\frac{a+c}{2}\qquad \textbf{(B)}\ b=\frac{a-c}{2}\qquad \textbf{(C)}\ a=\frac{b-c}{2} \qquad \textbf{(D)}\ a=\frac{b+c}{2}\qquad \textbf{(E)}\ \text{none of these}$


 

	\item (\textit{AHSME 1960}) To satisfy the equation $\frac{a+b}{a}=\frac{b}{a+b}$, $a$ and $b$ must be:


 $\textbf{(A)}\ \text{both rational}\qquad\textbf{(B)}\ \text{both real but not rational}\qquad\textbf{(C)}\ \text{both not real}\qquad$
$\textbf{(D)}\ \text{one real, one not real}\qquad\textbf{(E)}\ \text{one real, one not real or both not real}$


 



 

	\item (\textit{AHSME 1960}) Given right $\triangle ABC$ with legs $BC=3, AC=4$. Find the length of the shorter angle trisector from $C$ to the hypotenuse:


 $\textbf{(A)}\ \frac{32\sqrt{3}-24}{13}\qquad\textbf{(B)}\ \frac{12\sqrt{3}-9}{13}\qquad\textbf{(C)}\ 6\sqrt{3}-8\qquad\textbf{(D)}\ \frac{5\sqrt{10}}{6}\qquad\textbf{(E)}\ \frac{25}{12}\qquad$


 

\end{enumerate}